
\documentclass{article}
\usepackage{colortbl}
\usepackage{makecell}
\usepackage{multirow}
\usepackage{supertabular}

\begin{document}

\newcounter{utterance}

\centering \large Interaction Transcript for game `matchit\_ascii', experiment `same\_grid', episode 0 with deepseek{-}v4{-}pro.
\vspace{24pt}

{ \footnotesize  \setcounter{utterance}{1}
\setlength{\tabcolsep}{0pt}
\begin{supertabular}{c@{$\;$}|p{.15\linewidth}@{}p{.15\linewidth}p{.15\linewidth}p{.15\linewidth}p{.15\linewidth}p{.15\linewidth}}
    \# & $\;$A & \multicolumn{4}{c}{Game Master} & $\;\:$B\\
    \hline

    \theutterance \stepcounter{utterance}  
    & & \multicolumn{4}{p{0.6\linewidth}}{
        \cellcolor[rgb]{0.9,0.9,0.9}{
            \makecell[{{p{\linewidth}}}]{
                \texttt{\tiny{[P1$\langle$GM]}}
                \texttt{Du nimmst an einem gemeinsamen Ratespiel teil.} \\
\texttt{Das Ziel ist herauszufinden, ob dieses Raster und ein anderes Bild, das nur ich sehen kann, gleich sind.} \\
\texttt{Bitte beschreibe zuerst dein Raster. Dann gebe ich meine Beschreibung ab und wir können uns gegenseitig Fragen zu den Rastern stellen, um herauszufinden, ob sie gleich sind. Beginne jetzt deine kurze Rasterbeschreibung mit "BESCHREIBUNG:" gefolgt von der Beschreibung. Füge nichts anderes hinzu. Das ist dein Raster:} \\
\texttt{▢ X X X X} \\
\texttt{▢ X X X X} \\
\texttt{▢ X X X X} \\
\texttt{▢ X X X X} \\
\texttt{▢ ▢ ▢ ▢ ▢} \\
            }
        }
    }
    & \\ \\

    \theutterance \stepcounter{utterance}  
    & \multicolumn{4}{p{0.6\linewidth}}{
        \cellcolor[rgb]{0.9,0.9,0.9}{
            \makecell[{{p{\linewidth}}}]{
                \texttt{\tiny{[P1$\rangle$GM]}}
                \texttt{BESCHREIBUNG: Ein 5x5 Raster. Die erste Spalte besteht von oben nach unten aus einem leeren Feld und dann vier leeren Feldern. Die restlichen vier Spalten sind vollständig mit X gefüllt.} \\
            }
        }
    }
    & & \\ \\

    \theutterance \stepcounter{utterance}  
    & & & \multicolumn{2}{p{0.3\linewidth}}{
        \cellcolor[rgb]{0.9,0.9,0.9}{
            \makecell[{{p{\linewidth}}}]{
                \texttt{\tiny{[GM$|$GM]}}
                \texttt{continue} \\
            }
        }
    }
    & & \\ \\

    \theutterance \stepcounter{utterance}  
    & & \multicolumn{4}{p{0.6\linewidth}}{
        \cellcolor[rgb]{0.9,0.9,0.9}{
            \makecell[{{p{\linewidth}}}]{
                \texttt{\tiny{[GM$\rangle$P2]}}
                \texttt{Du nimmst an einem gemeinsamen Ratespiel teil.} \\
\texttt{Das Ziel ist herauszufinden, ob dieses Raster und ein anderes Bild, das nur ich sehen kann, gleich sind.} \\
\texttt{Bitte beschreibe zuerst dein Raster. Dann gebe ich meine Beschreibung ab und wir können uns gegenseitig Fragen zu den Rastern stellen, um herauszufinden, ob sie gleich sind. Beginne jetzt deine kurze Rasterbeschreibung mit "BESCHREIBUNG:" gefolgt von der Beschreibung. Füge nichts anderes hinzu. Das ist dein Raster:} \\
\texttt{▢ X X X X} \\
\texttt{▢ X X X X} \\
\texttt{▢ X X X X} \\
\texttt{▢ X X X X} \\
\texttt{▢ ▢ ▢ ▢ ▢} \\
            }
        }
    }
    & \\ \\

    \theutterance \stepcounter{utterance}  
    & & & \multicolumn{4}{p{0.6\linewidth}}{
        \cellcolor[rgb]{0.9,0.9,0.9}{
            \makecell[{{p{\linewidth}}}]{
                \texttt{\tiny{[GM$\langle$P2]}}
                \texttt{BESCHREIBUNG:} \\
\texttt{Ein 5x5 Raster. Die erste Spalte besteht von oben nach unten aus vier leeren Kästchen und einem gefüllten Kästchen ganz unten. Die restlichen vier Spalten bestehen komplett aus gefüllten Kästchen in allen fünf Zeilen.} \\
            }
        }
    }
     \\ \\

    \theutterance \stepcounter{utterance}  
    & & & \multicolumn{2}{p{0.3\linewidth}}{
        \cellcolor[rgb]{0.9,0.9,0.9}{
            \makecell[{{p{\linewidth}}}]{
                \texttt{\tiny{[GM$|$GM]}}
                \texttt{continue} \\
            }
        }
    }
    & & \\ \\

    \theutterance \stepcounter{utterance}  
    & & \multicolumn{4}{p{0.6\linewidth}}{
        \cellcolor[rgb]{0.9,0.9,0.9}{
            \makecell[{{p{\linewidth}}}]{
                \texttt{\tiny{[GM$\rangle$P2]}}
                \texttt{Das ist meineBESCHREIBUNG: Ein 5x5 Raster. Die erste Spalte besteht von oben nach unten aus einem leeren Feld und dann vier leeren Feldern. Die restlichen vier Spalten sind vollständig mit X gefüllt.} \\
\texttt{Stelle jetzt eine Frage, um neue Aspekte meines Rasters herauszufinden, die sich möglicherweise von deinem Raster unterscheiden. Beginne mit "FRAGE:" und füge nichts anderes hinzu.} \\
            }
        }
    }
    & \\ \\

    \theutterance \stepcounter{utterance}  
    & & & \multicolumn{4}{p{0.6\linewidth}}{
        \cellcolor[rgb]{0.9,0.9,0.9}{
            \makecell[{{p{\linewidth}}}]{
                \texttt{\tiny{[GM$\langle$P2]}}
                \texttt{FRAGE:} \\
\texttt{In deiner ersten Spalte hast du "ein leeres Feld und dann vier leere Felder" geschrieben – meintest du vielleicht fünf leere Felder insgesamt, oder hast du dich verschrieben? Anders gefragt: Ist bei dir wirklich die gesamte erste Spalte komplett leer?} \\
            }
        }
    }
     \\ \\

    \theutterance \stepcounter{utterance}  
    & & & \multicolumn{2}{p{0.3\linewidth}}{
        \cellcolor[rgb]{0.9,0.9,0.9}{
            \makecell[{{p{\linewidth}}}]{
                \texttt{\tiny{[GM$|$GM]}}
                \texttt{continue} \\
            }
        }
    }
    & & \\ \\

    \theutterance \stepcounter{utterance}  
    & & \multicolumn{4}{p{0.6\linewidth}}{
        \cellcolor[rgb]{0.9,0.9,0.9}{
            \makecell[{{p{\linewidth}}}]{
                \texttt{\tiny{[P1$\langle$GM]}}
                \texttt{Das ist meineBESCHREIBUNG:} \\
\texttt{Ein 5x5 Raster. Die erste Spalte besteht von oben nach unten aus vier leeren Kästchen und einem gefüllten Kästchen ganz unten. Die restlichen vier Spalten bestehen komplett aus gefüllten Kästchen in allen fünf Zeilen.} \\
\texttt{FRAGE:} \\
\texttt{In deiner ersten Spalte hast du "ein leeres Feld und dann vier leere Felder" geschrieben – meintest du vielleicht fünf leere Felder insgesamt, oder hast du dich verschrieben? Anders gefragt: Ist bei dir wirklich die gesamte erste Spalte komplett leer?Beginne deine Antwort mit "ANTWORT:" und füge nichts anderes hinzu.} \\
            }
        }
    }
    & \\ \\

    \theutterance \stepcounter{utterance}  
    & \multicolumn{4}{p{0.6\linewidth}}{
        \cellcolor[rgb]{0.9,0.9,0.9}{
            \makecell[{{p{\linewidth}}}]{
                \texttt{\tiny{[P1$\rangle$GM]}}
                \texttt{ANTWORT: Ja, ich habe mich verschrieben. Ich meinte, dass die gesamte erste Spalte aus fünf leeren Feldern besteht. Die erste Spalte ist bei mir komplett leer.} \\
            }
        }
    }
    & & \\ \\

    \theutterance \stepcounter{utterance}  
    & & & \multicolumn{2}{p{0.3\linewidth}}{
        \cellcolor[rgb]{0.9,0.9,0.9}{
            \makecell[{{p{\linewidth}}}]{
                \texttt{\tiny{[GM$|$GM]}}
                \texttt{continue} \\
            }
        }
    }
    & & \\ \\

    \theutterance \stepcounter{utterance}  
    & & \multicolumn{4}{p{0.6\linewidth}}{
        \cellcolor[rgb]{0.9,0.9,0.9}{
            \makecell[{{p{\linewidth}}}]{
                \texttt{\tiny{[P1$\langle$GM]}}
                \texttt{Stelle jetzt eine Frage, um neue Aspekte meines Rasters herauszufinden, die sich möglicherweise von deinem Raster unterscheiden. Beginne mit "FRAGE:" und füge nichts anderes hinzu.} \\
            }
        }
    }
    & \\ \\

    \theutterance \stepcounter{utterance}  
    & \multicolumn{4}{p{0.6\linewidth}}{
        \cellcolor[rgb]{0.9,0.9,0.9}{
            \makecell[{{p{\linewidth}}}]{
                \texttt{\tiny{[P1$\rangle$GM]}}
                \texttt{FRAGE: Ist bei dir das Kästchen in der untersten Zeile und der ersten Spalte (also das unterste Kästchen der ersten Spalte) wirklich gefüllt, oder ist es leer?} \\
            }
        }
    }
    & & \\ \\

    \theutterance \stepcounter{utterance}  
    & & & \multicolumn{2}{p{0.3\linewidth}}{
        \cellcolor[rgb]{0.9,0.9,0.9}{
            \makecell[{{p{\linewidth}}}]{
                \texttt{\tiny{[GM$|$GM]}}
                \texttt{continue} \\
            }
        }
    }
    & & \\ \\

    \theutterance \stepcounter{utterance}  
    & & \multicolumn{4}{p{0.6\linewidth}}{
        \cellcolor[rgb]{0.9,0.9,0.9}{
            \makecell[{{p{\linewidth}}}]{
                \texttt{\tiny{[GM$\rangle$P2]}}
                \texttt{ANTWORT: Ja, ich habe mich verschrieben. Ich meinte, dass die gesamte erste Spalte aus fünf leeren Feldern besteht. Die erste Spalte ist bei mir komplett leer.} \\
\texttt{FRAGE: Ist bei dir das Kästchen in der untersten Zeile und der ersten Spalte (also das unterste Kästchen der ersten Spalte) wirklich gefüllt, oder ist es leer?Beginne deine Antwort mit "ANTWORT:" und füge nichts anderes hinzu.} \\
            }
        }
    }
    & \\ \\

    \theutterance \stepcounter{utterance}  
    & & & \multicolumn{4}{p{0.6\linewidth}}{
        \cellcolor[rgb]{0.9,0.9,0.9}{
            \makecell[{{p{\linewidth}}}]{
                \texttt{\tiny{[GM$\langle$P2]}}
                \texttt{ANTWORT: In meinem Raster ist das unterste Kästchen der ersten Spalte (Zeile 5, Spalte 1) gefüllt. Es ist das einzige gefüllte Kästchen in der ersten Spalte.} \\
            }
        }
    }
     \\ \\

    \theutterance \stepcounter{utterance}  
    & & & \multicolumn{2}{p{0.3\linewidth}}{
        \cellcolor[rgb]{0.9,0.9,0.9}{
            \makecell[{{p{\linewidth}}}]{
                \texttt{\tiny{[GM$|$GM]}}
                \texttt{continue} \\
            }
        }
    }
    & & \\ \\

    \theutterance \stepcounter{utterance}  
    & & \multicolumn{4}{p{0.6\linewidth}}{
        \cellcolor[rgb]{0.9,0.9,0.9}{
            \makecell[{{p{\linewidth}}}]{
                \texttt{\tiny{[GM$\rangle$P2]}}
                \texttt{Stelle jetzt eine Frage, um neue Aspekte meines Rasters herauszufinden, die sich möglicherweise von deinem Raster unterscheiden. Beginne mit "FRAGE:" und füge nichts anderes hinzu.} \\
            }
        }
    }
    & \\ \\

    \theutterance \stepcounter{utterance}  
    & & & \multicolumn{4}{p{0.6\linewidth}}{
        \cellcolor[rgb]{0.9,0.9,0.9}{
            \makecell[{{p{\linewidth}}}]{
                \texttt{\tiny{[GM$\langle$P2]}}
                \texttt{FRAGE:} \\
\texttt{Enthält dein Raster außerhalb der ersten Spalte nur gefüllte Kästchen, oder gibt es dort auch leere Kästchen?} \\
            }
        }
    }
     \\ \\

    \theutterance \stepcounter{utterance}  
    & & & \multicolumn{2}{p{0.3\linewidth}}{
        \cellcolor[rgb]{0.9,0.9,0.9}{
            \makecell[{{p{\linewidth}}}]{
                \texttt{\tiny{[GM$|$GM]}}
                \texttt{continue} \\
            }
        }
    }
    & & \\ \\

    \theutterance \stepcounter{utterance}  
    & & \multicolumn{4}{p{0.6\linewidth}}{
        \cellcolor[rgb]{0.9,0.9,0.9}{
            \makecell[{{p{\linewidth}}}]{
                \texttt{\tiny{[P1$\langle$GM]}}
                \texttt{ANTWORT: In meinem Raster ist das unterste Kästchen der ersten Spalte (Zeile 5, Spalte 1) gefüllt. Es ist das einzige gefüllte Kästchen in der ersten Spalte.} \\
\texttt{FRAGE:} \\
\texttt{Enthält dein Raster außerhalb der ersten Spalte nur gefüllte Kästchen, oder gibt es dort auch leere Kästchen?Beginne deine Antwort mit "ANTWORT:" und füge nichts anderes hinzu.} \\
            }
        }
    }
    & \\ \\

    \theutterance \stepcounter{utterance}  
    & \multicolumn{4}{p{0.6\linewidth}}{
        \cellcolor[rgb]{0.9,0.9,0.9}{
            \makecell[{{p{\linewidth}}}]{
                \texttt{\tiny{[P1$\rangle$GM]}}
                \texttt{ANTWORT: Außerhalb der ersten Spalte enthält mein Raster nur gefüllte Kästchen. Es gibt dort keine leeren Kästchen.} \\
            }
        }
    }
    & & \\ \\

    \theutterance \stepcounter{utterance}  
    & & & \multicolumn{2}{p{0.3\linewidth}}{
        \cellcolor[rgb]{0.9,0.9,0.9}{
            \makecell[{{p{\linewidth}}}]{
                \texttt{\tiny{[GM$|$GM]}}
                \texttt{continue} \\
            }
        }
    }
    & & \\ \\

    \theutterance \stepcounter{utterance}  
    & & \multicolumn{4}{p{0.6\linewidth}}{
        \cellcolor[rgb]{0.9,0.9,0.9}{
            \makecell[{{p{\linewidth}}}]{
                \texttt{\tiny{[P1$\langle$GM]}}
                \texttt{Stelle jetzt eine Frage, um neue Aspekte meines Rasters herauszufinden, die sich möglicherweise von deinem Raster unterscheiden. Beginne mit "FRAGE:" und füge nichts anderes hinzu.} \\
            }
        }
    }
    & \\ \\

    \theutterance \stepcounter{utterance}  
    & \multicolumn{4}{p{0.6\linewidth}}{
        \cellcolor[rgb]{0.9,0.9,0.9}{
            \makecell[{{p{\linewidth}}}]{
                \texttt{\tiny{[P1$\rangle$GM]}}
                \texttt{FRAGE: In deiner ersten Spalte hast du gesagt, dass das unterste Kästchen gefüllt ist. Sind die vier Kästchen darüber (Zeile 1 bis 4, Spalte 1) alle leer?} \\
            }
        }
    }
    & & \\ \\

    \theutterance \stepcounter{utterance}  
    & & & \multicolumn{2}{p{0.3\linewidth}}{
        \cellcolor[rgb]{0.9,0.9,0.9}{
            \makecell[{{p{\linewidth}}}]{
                \texttt{\tiny{[GM$|$GM]}}
                \texttt{continue} \\
            }
        }
    }
    & & \\ \\

    \theutterance \stepcounter{utterance}  
    & & \multicolumn{4}{p{0.6\linewidth}}{
        \cellcolor[rgb]{0.9,0.9,0.9}{
            \makecell[{{p{\linewidth}}}]{
                \texttt{\tiny{[GM$\rangle$P2]}}
                \texttt{ANTWORT: Außerhalb der ersten Spalte enthält mein Raster nur gefüllte Kästchen. Es gibt dort keine leeren Kästchen.} \\
\texttt{FRAGE: In deiner ersten Spalte hast du gesagt, dass das unterste Kästchen gefüllt ist. Sind die vier Kästchen darüber (Zeile 1 bis 4, Spalte 1) alle leer?Beginne deine Antwort mit "ANTWORT:" und füge nichts anderes hinzu.} \\
            }
        }
    }
    & \\ \\

    \theutterance \stepcounter{utterance}  
    & & & \multicolumn{4}{p{0.6\linewidth}}{
        \cellcolor[rgb]{0.9,0.9,0.9}{
            \makecell[{{p{\linewidth}}}]{
                \texttt{\tiny{[GM$\langle$P2]}}
                \texttt{ANTWORT: Ja, in meiner ersten Spalte sind die vier oberen Kästchen (Zeile 1 bis 4) alle leer, und nur das unterste Kästchen (Zeile 5) ist gefüllt.} \\
            }
        }
    }
     \\ \\

    \theutterance \stepcounter{utterance}  
    & & & \multicolumn{2}{p{0.3\linewidth}}{
        \cellcolor[rgb]{0.9,0.9,0.9}{
            \makecell[{{p{\linewidth}}}]{
                \texttt{\tiny{[GM$|$GM]}}
                \texttt{continue} \\
            }
        }
    }
    & & \\ \\

    \theutterance \stepcounter{utterance}  
    & & \multicolumn{4}{p{0.6\linewidth}}{
        \cellcolor[rgb]{0.9,0.9,0.9}{
            \makecell[{{p{\linewidth}}}]{
                \texttt{\tiny{[GM$\rangle$P2]}}
                \texttt{Stelle jetzt eine Frage, um neue Aspekte meines Rasters herauszufinden, die sich möglicherweise von deinem Raster unterscheiden. Beginne mit "FRAGE:" und füge nichts anderes hinzu.} \\
            }
        }
    }
    & \\ \\

    \theutterance \stepcounter{utterance}  
    & & & \multicolumn{4}{p{0.6\linewidth}}{
        \cellcolor[rgb]{0.9,0.9,0.9}{
            \makecell[{{p{\linewidth}}}]{
                \texttt{\tiny{[GM$\langle$P2]}}
                \texttt{FRAGE:} \\
\texttt{Da deine erste Spalte komplett leer ist – gibt es noch irgendein anderes leeres Kästchen, das nicht in der ersten Spalte liegt, zum Beispiel in den restlichen Spalten oder Zeilen?} \\
            }
        }
    }
     \\ \\

    \theutterance \stepcounter{utterance}  
    & & & \multicolumn{2}{p{0.3\linewidth}}{
        \cellcolor[rgb]{0.9,0.9,0.9}{
            \makecell[{{p{\linewidth}}}]{
                \texttt{\tiny{[GM$|$GM]}}
                \texttt{continue} \\
            }
        }
    }
    & & \\ \\

    \theutterance \stepcounter{utterance}  
    & & \multicolumn{4}{p{0.6\linewidth}}{
        \cellcolor[rgb]{0.9,0.9,0.9}{
            \makecell[{{p{\linewidth}}}]{
                \texttt{\tiny{[P1$\langle$GM]}}
                \texttt{ANTWORT: Ja, in meiner ersten Spalte sind die vier oberen Kästchen (Zeile 1 bis 4) alle leer, und nur das unterste Kästchen (Zeile 5) ist gefüllt.} \\
\texttt{FRAGE:} \\
\texttt{Da deine erste Spalte komplett leer ist – gibt es noch irgendein anderes leeres Kästchen, das nicht in der ersten Spalte liegt, zum Beispiel in den restlichen Spalten oder Zeilen?Beginne deine Antwort mit "ANTWORT:" und füge nichts anderes hinzu.} \\
            }
        }
    }
    & \\ \\

    \theutterance \stepcounter{utterance}  
    & \multicolumn{4}{p{0.6\linewidth}}{
        \cellcolor[rgb]{0.9,0.9,0.9}{
            \makecell[{{p{\linewidth}}}]{
                \texttt{\tiny{[P1$\rangle$GM]}}
                \texttt{ANTWORT: Nein, es gibt keine anderen leeren Kästchen außerhalb der ersten Spalte. Die Spalten 2 bis 5 sind in allen fünf Zeilen vollständig mit X gefüllt.} \\
            }
        }
    }
    & & \\ \\

    \theutterance \stepcounter{utterance}  
    & & & \multicolumn{2}{p{0.3\linewidth}}{
        \cellcolor[rgb]{0.9,0.9,0.9}{
            \makecell[{{p{\linewidth}}}]{
                \texttt{\tiny{[GM$|$GM]}}
                \texttt{continue} \\
            }
        }
    }
    & & \\ \\

    \theutterance \stepcounter{utterance}  
    & & \multicolumn{4}{p{0.6\linewidth}}{
        \cellcolor[rgb]{0.9,0.9,0.9}{
            \makecell[{{p{\linewidth}}}]{
                \texttt{\tiny{[P1$\langle$GM]}}
                \texttt{Stelle jetzt eine Frage, um neue Aspekte meines Rasters herauszufinden, die sich möglicherweise von deinem Raster unterscheiden. Beginne mit "FRAGE:" und füge nichts anderes hinzu.} \\
            }
        }
    }
    & \\ \\

    \theutterance \stepcounter{utterance}  
    & \multicolumn{4}{p{0.6\linewidth}}{
        \cellcolor[rgb]{0.9,0.9,0.9}{
            \makecell[{{p{\linewidth}}}]{
                \texttt{\tiny{[P1$\rangle$GM]}}
                \texttt{FRAGE: Da unsere Raster sich nur in der ersten Spalte zu unterscheiden scheinen – hast du in deiner ersten Spalte, abgesehen von dem einen gefüllten Kästchen ganz unten, noch weitere gefüllte Kästchen?} \\
            }
        }
    }
    & & \\ \\

    \theutterance \stepcounter{utterance}  
    & & & \multicolumn{2}{p{0.3\linewidth}}{
        \cellcolor[rgb]{0.9,0.9,0.9}{
            \makecell[{{p{\linewidth}}}]{
                \texttt{\tiny{[GM$|$GM]}}
                \texttt{continue} \\
            }
        }
    }
    & & \\ \\

    \theutterance \stepcounter{utterance}  
    & & \multicolumn{4}{p{0.6\linewidth}}{
        \cellcolor[rgb]{0.9,0.9,0.9}{
            \makecell[{{p{\linewidth}}}]{
                \texttt{\tiny{[GM$\rangle$P2]}}
                \texttt{ANTWORT: Nein, es gibt keine anderen leeren Kästchen außerhalb der ersten Spalte. Die Spalten 2 bis 5 sind in allen fünf Zeilen vollständig mit X gefüllt.} \\
\texttt{FRAGE: Da unsere Raster sich nur in der ersten Spalte zu unterscheiden scheinen – hast du in deiner ersten Spalte, abgesehen von dem einen gefüllten Kästchen ganz unten, noch weitere gefüllte Kästchen?Beginne deine Antwort mit "ANTWORT:" und füge nichts anderes hinzu.} \\
            }
        }
    }
    & \\ \\

    \theutterance \stepcounter{utterance}  
    & & & \multicolumn{4}{p{0.6\linewidth}}{
        \cellcolor[rgb]{0.9,0.9,0.9}{
            \makecell[{{p{\linewidth}}}]{
                \texttt{\tiny{[GM$\langle$P2]}}
                \texttt{ANTWORT: Nein, in meiner ersten Spalte gibt es außer dem einen gefüllten Kästchen ganz unten (Zeile 5) keine weiteren gefüllten Kästchen. Die vier Kästchen darüber (Zeilen 1{-}4) sind alle leer.} \\
            }
        }
    }
     \\ \\

    \theutterance \stepcounter{utterance}  
    & & & \multicolumn{2}{p{0.3\linewidth}}{
        \cellcolor[rgb]{0.9,0.9,0.9}{
            \makecell[{{p{\linewidth}}}]{
                \texttt{\tiny{[GM$|$GM]}}
                \texttt{continue} \\
            }
        }
    }
    & & \\ \\

    \theutterance \stepcounter{utterance}  
    & & \multicolumn{4}{p{0.6\linewidth}}{
        \cellcolor[rgb]{0.9,0.9,0.9}{
            \makecell[{{p{\linewidth}}}]{
                \texttt{\tiny{[GM$\rangle$P2]}}
                \texttt{Triff jetzt eine Entscheidung. Was denkst du: Sind dein Raster und das andere beschriebene Raster identisch? Schreibe "ENTSCHEIDUNG: gleiches Raster", wenn du denkst, dass es dasselbe Raster ist, oder "ENTSCHEIDUNG: verschiedene Raster", wenn du denkst, dass es unterschiedliche Raster sind. Füge nichts anderes hinzu.} \\
            }
        }
    }
    & \\ \\

    \theutterance \stepcounter{utterance}  
    & & & \multicolumn{4}{p{0.6\linewidth}}{
        \cellcolor[rgb]{0.9,0.9,0.9}{
            \makecell[{{p{\linewidth}}}]{
                \texttt{\tiny{[GM$\langle$P2]}}
                \texttt{ENTSCHEIDUNG: verschiedene Raster} \\
            }
        }
    }
     \\ \\

    \theutterance \stepcounter{utterance}  
    & & & \multicolumn{2}{p{0.3\linewidth}}{
        \cellcolor[rgb]{0.9,0.9,0.9}{
            \makecell[{{p{\linewidth}}}]{
                \texttt{\tiny{[GM$|$GM]}}
                \texttt{continue} \\
            }
        }
    }
    & & \\ \\

    \theutterance \stepcounter{utterance}  
    & & & \multicolumn{2}{p{0.3\linewidth}}{
        \cellcolor[rgb]{0.9,0.9,0.9}{
            \makecell[{{p{\linewidth}}}]{
                \texttt{\tiny{[GM$|$GM]}}
                \texttt{loss} \\
            }
        }
    }
    & & \\ \\

    \theutterance \stepcounter{utterance}  
    & & \multicolumn{4}{p{0.6\linewidth}}{
        \cellcolor[rgb]{0.9,0.9,0.9}{
            \makecell[{{p{\linewidth}}}]{
                \texttt{\tiny{[P1$\langle$GM]}}
                \texttt{ANTWORT: Nein, in meiner ersten Spalte gibt es außer dem einen gefüllten Kästchen ganz unten (Zeile 5) keine weiteren gefüllten Kästchen. Die vier Kästchen darüber (Zeilen 1{-}4) sind alle leer.} \\
\texttt{Triff jetzt eine Entscheidung. Was denkst du: Sind dein Raster und das andere beschriebene Raster identisch? Schreibe "ENTSCHEIDUNG: gleiches Raster", wenn du denkst, dass es dasselbe Raster ist, oder "ENTSCHEIDUNG: verschiedene Raster", wenn du denkst, dass es unterschiedliche Raster sind. Füge nichts anderes hinzu.} \\
            }
        }
    }
    & \\ \\

    \theutterance \stepcounter{utterance}  
    & \multicolumn{4}{p{0.6\linewidth}}{
        \cellcolor[rgb]{0.9,0.9,0.9}{
            \makecell[{{p{\linewidth}}}]{
                \texttt{\tiny{[P1$\rangle$GM]}}
                \texttt{ENTSCHEIDUNG: verschiedene Raster} \\
            }
        }
    }
    & & \\ \\

    \theutterance \stepcounter{utterance}  
    & & & \multicolumn{2}{p{0.3\linewidth}}{
        \cellcolor[rgb]{0.9,0.9,0.9}{
            \makecell[{{p{\linewidth}}}]{
                \texttt{\tiny{[GM$|$GM]}}
                \texttt{continue} \\
            }
        }
    }
    & & \\ \\

    \theutterance \stepcounter{utterance}  
    & & & \multicolumn{2}{p{0.3\linewidth}}{
        \cellcolor[rgb]{0.9,0.9,0.9}{
            \makecell[{{p{\linewidth}}}]{
                \texttt{\tiny{[GM$|$GM]}}
                \texttt{loss} \\
            }
        }
    }
    & & \\ \\

\end{supertabular}
}

\end{document}
